\documentclass[10pt,letterpaper]{article}
\usepackage[top=0.68in,bottom=0.70in,left=0.82in,right=0.82in]{geometry}
\usepackage{fontspec}
\IfFontExistsTF{Helvetica Neue}{
  \setmainfont{Helvetica Neue}[BoldFont=Helvetica Neue Medium,ItalicFont=Helvetica Neue Italic]
}{
  \setmainfont{texgyreheros-regular.otf}[BoldFont=texgyreheros-bold.otf,ItalicFont=texgyreheros-italic.otf]
}
\setmonofont{lmmono10-regular.otf}[Scale=0.92]
\usepackage{xcolor,enumitem,fancyhdr,tabularx,array}
\usepackage[unicode,pdfauthor={},pdftitle={Predict frontier difficulty across a task portfolio},hidelinks]{hyperref}
\definecolor{ink}{HTML}{1D1D1F}
\definecolor{secondary}{HTML}{626269}
\definecolor{hairline}{HTML}{DADADD}
\definecolor{surface}{HTML}{F5F5F7}
\definecolor{accent}{HTML}{0066CC}
\color{ink}
\setlength{\parindent}{0pt}
\setlength{\parskip}{0pt}
\renewcommand{\baselinestretch}{1.15}
\setlength{\emergencystretch}{2em}
\hyphenpenalty=10000
\setlist[itemize]{leftmargin=12pt,label={\color{secondary}\textbullet},itemsep=5pt,topsep=5pt,parsep=0pt}
\pagestyle{fancy}\fancyhf{}
\setlength{\footskip}{26pt}
\fancyfoot[R]{\fontsize{8}{10}\selectfont\color{secondary}\thepage\enspace /\enspace 2}
\renewcommand{\headrulewidth}{0pt}
\newcommand{\eyebrow}[1]{{\fontsize{8.5}{11}\selectfont\color{secondary}\MakeUppercase{#1}}\par}
\newcommand{\stage}[2]{\vspace{18pt}{\fontsize{17}{21}\selectfont\bfseries #2}\par\vspace{9pt}}
\newcommand{\subhead}[1]{{\fontsize{10.5}{13}\selectfont\bfseries #1}\par\vspace{5pt}}
\newcommand{\ruleline}{\vspace{15pt}{\color{hairline}\hrule height 0.4pt}\vspace{0pt}}
\newcommand{\question}[3]{\noindent\begin{minipage}[t]{0.065\linewidth}\vspace{0pt}{\fontsize{9}{12}\selectfont\color{secondary}#1}\end{minipage}\begin{minipage}[t]{0.935\linewidth}\vspace{0pt}\subhead{#2}#3\end{minipage}\par\vspace{13pt}}
\begin{document}
\raggedright
{\fontsize{29}{33}\selectfont\bfseries Predict frontier difficulty\\across a task portfolio.\par}
\vspace{13pt}
Build up to five agent tasks in any domain, aiming for at least one that is difficult for a frontier model. Success requires at least one valid task and a defensible forecast made before frontier results are known.
\vspace{13pt}

\begingroup
\setlength{\fboxsep}{12pt}
\colorbox{surface}{\begin{minipage}{\dimexpr\linewidth-24pt\relax}
\begin{tabularx}{\linewidth}{@{}X X X@{}}
{\fontsize{20}{24}\selectfont\bfseries 5} & {\fontsize{20}{24}\selectfont\bfseries 8} & {\fontsize{20}{24}\selectfont\bfseries 0--2}\\[3pt]
{\color{secondary}Final task packages} & {\color{secondary}Independent trials per task} & {\color{secondary}Frontier passes to be hard}
\end{tabularx}
\end{minipage}}
\endgroup
\vspace{11pt}

Evaluators use either Fable or GPT-5.6 Sol. \textbf{Do not run either target model yourself.} A hard task must have a passing canonical solution and a fair verifier. Infrastructure and harness failures do not count as model failures.

\ruleline
\stage{DAY 1}{Build and preregister}
Submit five separate packages in Harbor format. Tasks may share a base repository or environment; each needs its own instruction, verifier, and canonical solution.
\vspace{12pt}

\begin{minipage}[t]{0.51\linewidth}
\vspace{0pt}\raggedright
\subhead{Package structure}
{\fontsize{9}{13}\selectfont\ttfamily
<task-slug>/\\
\hspace*{10pt}task.toml\\
\hspace*{10pt}instruction.md\\
\hspace*{10pt}environment/Dockerfile\\
\hspace*{10pt}tests/test.sh\\
\hspace*{10pt}solution/solve.sh\\
\hspace*{10pt}solution/solution.patch}\par
\vspace{5pt}
{\fontsize{9}{12}\selectfont\color{secondary}The patch may be replaced by another canonical reference artifact.}
\end{minipage}\hfill
\begin{minipage}[t]{0.45\linewidth}
\vspace{0pt}\raggedright
\subhead{Validate every task}
\begin{itemize}[topsep=0pt,itemsep=7pt]
\item The canonical solution receives full reward.
\item A plausible incorrect or shortcut solution fails.
\item The verifier accepts materially equivalent correct solutions.
\end{itemize}
\end{minipage}\par
\vspace{16pt}

\subhead{Preregistration\enspace /\enspace \texttt{day1.md}}
Include one compact row per task covering:
\begin{itemize}[itemsep=4pt]
\item The capability gap and concrete failure mechanism.
\item The predicted number of frontier passes out of eight.
\item The evidence supporting that prediction.
\item The prompts, agents, tools, and orchestration used to create and critique it.
\end{itemize}
\vspace{5pt}
Rank all five tasks from most to least likely to be hard.

\newpage
{\fontsize{29}{33}\selectfont\bfseries Update from the evidence.\par}
\vspace{14pt}
\subhead{Overnight\enspace /\enspace Proxy panel}
Evaluators validate the packages and run the same inexpensive model panel on all five tasks. The next day, they return per-trial rewards and complete trajectories. Evaluation runtime is not part of your working time.

\ruleline
\stage{DAY 2}{Read evidence and iterate}
Decide which tasks to repair, harden, leave unchanged, or retire. Submit the final five task packages and a one-page \texttt{report.md} covering these six questions.
\vspace{15pt}

\question{01}{Changes and rationale}{What did you change after seeing the proxy results, and why?}
\question{02}{Pass-rate analysis}{What do the proxy pass rates imply, accounting for the small number of trials?}
\question{03}{Trajectory evidence}{What step-indexed evidence supports the failure mechanism you believe will transfer to the frontier model?}
\question{04}{Final forecasts}{For each task, provide $\mathsf{P(K=0),\,P(K=1),\,\ldots,\,P(K=8)}$, where $\mathsf{K}$ is its frontier pass count across eight trials.}
\question{05}{Hardest task and falsification}{Which task do you expect to be hardest, and what result would falsify that belief?}
\question{06}{Agent workflow}{How did your workflow divide generation, implementation, critique, testing, and revision?}

{\color{hairline}\hrule height 0.4pt}
\vspace{14pt}
\subhead{Final evaluation}
Evaluators run eight independent trials per final task on either Fable or GPT-5.6 Sol. A task counts as hard when its canonical solution passes, its verifier is fair, its runs contain real agent work, and the frontier model passes no more than two of eight trials.\par
\vspace{12pt}
{\color{secondary}The main deliverable is evidence of a process: generate a portfolio, identify likely hard tasks before frontier evaluation, update rationally from cheaper trials, and concentrate effort until at least one task is both hard and valid.}
\end{document}
